% ----------------------------------------------------------------------
% Resumo em Português
%
% Este arquivo contém o resumo do trabalho em português. O título
% “Resumo” deve ser centralizado e não numerado. A
% primeira linha do texto do resumo deve iniciar logo abaixo do título
% sem recuo. Usar espaço simples neste trecho conforme as normas.

% O título do resumo aparece centralizado, com mesma fonte das listas, e não é
% incluído no sumário.
{\centering
\large\bfseries Resumo\par}
% Insere 1,5cm de espaço vertical entre o título e o texto do resumo【543391667091976†L423-L441】.
\vspace{1.5cm}

% O resumo deve ser digitado em espaço simples.  Usa-se o
% ambiente singlespace fornecido pelo pacote setspace para aplicar
% espaçamento simples apenas nesta seção.
\begin{singlespace}
A duração gestacional termina em dois eventos mutuamente exclusivos, óbito
fetal e nascimento vivo, que frequentemente são analisados separadamente
apesar de sua natureza competitiva. Esta dissertação avaliou o tempo até o
parto no estado de São Paulo, combinando métodos de riscos competitivos com
uma estrutura bayesiana de fragilidade espacial compartilhada no nível das
Redes Regionais de Atenção à Saúde (RRAS). Foram utilizados microdados
públicos de 2023 do SIM-DOFET e do SINASC, harmonizados e analisados após
filtragem por casos completos, resultando em 3.266 óbitos fetais e 485.922
nascimentos vivos. As análises descritivas e as funções de incidência
acumulada mostraram forte assimetria entre os desfechos: os óbitos fetais se
concentraram em idades gestacionais mais precoces, enquanto os nascimentos
vivos se concentraram próximo ao termo. As CIF estratificadas indicaram
padrões sociais e obstétricos relevantes, incluindo perfis menos favoráveis
nos estratos de menor escolaridade materna e em grupos com histórico
obstétrico desfavorável. Os indicadores por RRAS revelaram heterogeneidade
territorial marcada, com descompasso entre carga absoluta e indicadores
relativos, sustentando a modelagem explícita da vulnerabilidade regional
latente. O modelo
bayesiano proposto de riscos competitivos com fragilidade ICAR foi desenvolvido
para fornecer estimativas regionais interpretáveis e com quantificação de
incerteza, apoiando o planejamento em saúde materna e perinatal no Sistema
Único de Saúde (SUS).


\vspace{0.5cm}
\noindent\textbf{Palavras‑chave:} óbito fetal. nascimento vivo. riscos
competitivos. fragilidade espacial. Redes Regionais de Atenção à Saúde.
\end{singlespace}
